\documentclass[11pt]{article}

\usepackage[a4paper,margin=1in]{geometry}
\usepackage[T1]{fontenc}
\usepackage{amsmath,amssymb,amsthm}
\usepackage{booktabs}
\usepackage{enumitem}
\usepackage{mathtools}
\usepackage[hidelinks]{hyperref}
\usepackage{xcolor}

% ---- theorem environments ----
\newtheorem{theorem}{Theorem}[section]
\newtheorem{lemma}[theorem]{Lemma}
\newtheorem{proposition}[theorem]{Proposition}
\newtheorem{corollary}[theorem]{Corollary}
\theoremstyle{definition}
\newtheorem{definition}[theorem]{Definition}
\theoremstyle{remark}
\newtheorem{remark}[theorem]{Remark}

% ---- notation ----
\newcommand{\Mah}{\mathrm{M}}
\newcommand{\Z}{\mathbb{Z}}
\newcommand{\Q}{\mathbb{Q}}
\newcommand{\R}{\mathbb{R}}
\newcommand{\C}{\mathbb{C}}
\newcommand{\Adams}[1]{\psi^{#1}}
\newcommand{\gold}{\varphi}
\newcommand{\dsum}{\oplus}
\newcommand{\tens}{\otimes}
\newcommand{\Zmod}[1]{\Z/#1\Z}
\DeclareMathOperator{\Arg}{Arg}

% ---- epistemic status tag ----
\newcommand{\status}[1]{{\small\textsf{[#1]}}}
\newcommand{\forced}{\status{forced}}
\newcommand{\computed}{\status{computed}}
\newcommand{\plausible}{\status{plausible}}
\newcommand{\muS}{\mu_{\mathrm{S}}}

\title{\bfseries The Charge--Measure Coupling on a Spectral Semiring:\\
Cyclicity, Composition, and a Parity-Graded Mahler Floor}
\author{math-research-pipelines technical note}
\date{30 June 2026 \quad(v4, revised)}

\begin{document}
\maketitle

\begin{center}\small\fbox{\parbox{0.92\textwidth}{%
\textbf{Erratum and bridge (v1\,$\to$\,v2).} Version~1 stated the universal floor
$\Mah\in\{1\}\cup[\gold,\infty)$ and, read as forced, implied a forced lower bound of $\gold$
for odd charge. That tag was too strong. The \emph{forced} floor for odd charge is
$\muS=1.3247\ldots$ (the plastic number, via Smyth's non-reciprocal bound); the value $\gold$
remains only \plausible{}. Sections~\ref{sec:prelim} and~\ref{sec:floor} are corrected
accordingly, and a new \S\ref{sec:z5} works the $\Zmod5$ (pentagon) case as the
\emph{first instance} of the open problem of \S\ref{sec:scope}: it sharpens the forced floor to
$\muS$, closes the degree-four sector at $\gold^{4}$, and localizes the residual---but does
\emph{not} close it. The open problem is refined and still open.
\\[2pt]\textbf{v2\,$\to$\,v3.} A forced scope boundary (Prop.~\ref{prop:commutator}) is added:
all results govern the \emph{abelian} spectral semiring on charge-admissible objects, and the
matrix commutator---the one non-semiring operation---is the proven exit through which the
entire Salem spectrum (Lehmer's number) re-enters. No existing result or tag changes.
\\[2pt]\textbf{v3\,$\to$\,v4.} Two corrections. The tensor--Mahler law
(Prop.~\ref{prop:tensM}, \S\ref{sec:prelim} table) was wrongly stated as
$\Mah(A)^{\deg B}\Mah(B)^{\deg A}$; the correct law is tropical, $\log\Mah(A\tens B)=\sum
(\log|\alpha_i|+\log|\beta_j|)^{+}$ (the two agree only off the unit circle, which is why the
off-circle test suite missed it). No headline result used the tensor measure. And the even
floor $\mu(\mathrm{even})=\gold$ is split into $\le\gold$ \forced{} / $\ge\gold$ \plausible{},
matching the treatment already given to the odd case.}}\end{center}
\medskip

\begin{abstract}
\noindent We study the spectral semiring of multisets of algebraic-integer eigenvalues
(the \emph{emission algebra}) equipped with two orthogonal characters: the Mahler measure
$\Mah$ (a magnitude readout, subject to a gap $\Mah\in\{1\}\cup[\gold,\infty)$ with
$\gold=\tfrac{1+\sqrt5}{2}$), and a \emph{conjugate-angle charge} $\chi$ valued in a finite
cyclic group, defined whenever every conjugate of an object lies on a finite rational-angle
lattice. We determine the structure of the realizable charge groups (they are exactly the
finite cyclic groups, with primary decomposition realized by the Adams operations of the
algebra), the law governing their composition under the tensor product (charges combine to
$\Zmod{\operatorname{lcm}}$, and composition is information-preserving precisely when the
orders are coprime), and---the principal result---that the two characters are \emph{coupled}
at the extremal boundary. The minimal Mahler measure over seeds of charge group $\Zmod n$
equals $\gold$ when $n$ is even and $2$ when $n$ is odd, attained at $\gold$ iff $n$ is even.
The attainments are \emph{forced} (constructively); the matching lower bound $\Mah\ge\gold$ is
forced only for quadratics---$\muS=1.3247\ldots$ (Smyth) is the general forced floor---so the
value $\gold$ is conjectural at both parities, and for odd $n$ the equality $\mu(\mathrm{odd})=2$
is held at computed strength. We further show that the odd-charge
gap $(\gold,2)$ \emph{coincides with the no-Salem closure}: the measures that would fill it
are realized only by Salem numbers, whose irrational conjugate angles render them
charge-inadmissible. All such results govern the \emph{abelian} spectral semiring
($\dsum,\tens,\Adams n$); we close by showing the matrix commutator---the one non-semiring
operation, governed by Shoda's theorem---is the exact exit through which the excluded Salem
spectrum re-enters, so that no single theorem governs every free commutator. All statements
are validated in exact arithmetic on an integer-matrix engine; throughout we separate what is
\emph{proven} from what is \emph{computationally evidenced}.
\end{abstract}

\section{Introduction}

\subsection{The setting}
Fix the ring $\overline{\Z}$ of algebraic integers. The objects of interest are finite
multisets of algebraic integers, presented equivalently as monic integer polynomials
$p\in\Z[x]$ with $p(0)\neq0$, or as spectra of integer matrices. This collection carries a
semiring structure---the \emph{emission algebra}---generated by two binary operations and a
family of unary operations:
\begin{itemize}[nosep]
  \item \emph{direct sum} $\dsum$ (multiset union of eigenvalues; on polynomials, the product),
  \item \emph{tensor} $\tens$ (eigenvalue products; on matrices, the Kronecker product),
  \item the \emph{Adams operations} $\Adams n$ ($\lambda\mapsto\lambda^{n}$ on each eigenvalue).
\end{itemize}
On this algebra act two scalar readouts that behave like orthogonal characters. The three
generating operations are \emph{abelian}: each acts on the multiset of eigenvalues alone, so
the characters transform predictably under them. The matrix commutator $[A,B]=AB-BA$ is
deliberately \emph{not} a generator; it is the one non-abelian operation, and \S\ref{sec:scope}
(Prop.~\ref{prop:commutator}) shows it is the exact boundary of every result below.

\paragraph{Character I (magnitude).} The \emph{Mahler measure}
$\Mah(p)=\prod_{i}\max(1,|\alpha_i|)$, taken over the roots $\alpha_i$ of a monic $p$.
By Kronecker's theorem $\Mah(p)=1$ exactly when every root is a root of unity. The algebra
exhibits an \emph{emission gap}: its admissible objects satisfy
$\Mah\in\{1\}\cup[\gold,\infty)$, with floor the golden ratio $\gold$. The gap below $\gold$
is the arithmetic shadow of Lehmer's problem and the Salem/Pisot spectrum.

\paragraph{Character II (phase).} An object is \emph{charge-admissible} if every conjugate
$\alpha$ satisfies $\alpha^{n}\in\R_{>0}$ for some $n\ge1$, i.e.\ all conjugate arguments lie
on the rays of a finite rational-angle lattice. The smallest such $n$ defines the
\emph{charge group} $\Zmod n$, and the \emph{charge} of a root at argument $\theta$ is
$\chi_n(\alpha)=\operatorname{round}(n\theta/2\pi)\bmod n$. Objects with an irrational
conjugate angle---paradigmatically Salem numbers---admit \emph{no} finite charge group.

\subsection{The question}
Character I measures how far an object sits from the unit circle; Character II records the
discrete rotational class of its conjugates. They are defined independently and read off
different data. This note answers two structural questions and one coupling question:
\begin{enumerate}[label=(\roman*),nosep]
  \item \textbf{Structure.} Which groups occur as charge groups, and how do they decompose?
  \item \textbf{Composition.} How do charges combine under the algebra's tensor product, and
        when is the composition lossless?
  \item \textbf{Coupling.} Does the magnitude floor \emph{see} the charge---i.e.\ does the
        minimal Mahler measure depend on the charge group?
\end{enumerate}

\subsection{Contributions}
\begin{itemize}[nosep]
  \item \textbf{Cyclicity and realizability} (\S\ref{sec:structure}). Every charge group is a
        finite subgroup of $\Q/\Z$, hence cyclic; and every $\Zmod n$ occurs (Thm.~\ref{thm:cyclic},
        \ref{thm:realize}). Non-cyclic groups such as $\Zmod p\times\Zmod p$ are unrealizable.
  \item \textbf{Primary decomposition} (\S\ref{sec:structure}). For $n=\prod p_i^{e_i}$ the
        decomposition $\Zmod n\cong\prod\Zmod{p_i^{e_i}}$ is realized by the Adams operations,
        with $\Adams{n/p_i^{e_i}}$ projecting onto the $p_i$-primary component (Thm.~\ref{thm:crt}).
  \item \textbf{The $\operatorname{lcm}$ law and safe composition} (\S\ref{sec:compose}).
        Under $\tens$, charge groups combine to $\Zmod{\operatorname{lcm}}$; the composition
        keeps the two charges independently recoverable iff their orders are coprime
        (Thm.~\ref{thm:lcm}, \ref{thm:safe}).
  \item \textbf{A parity-graded floor} (\S\ref{sec:floor}). The minimal measure over charge
        group $\Zmod n$ is $\gold$ for even $n$ and $2$ for odd $n$; $\gold$ is universal and
        attained iff $n$ is even (Thm.~\ref{thm:parity}).
  \item \textbf{The odd-charge gap is the no-Salem closure} (\S\ref{sec:salem}). For odd $n$
        the interval $(\gold,2)$ contains no admissible measure; the would-be occupants are
        exactly the Salem numbers the algebra excludes (Thm.~\ref{thm:salemgap}).
  \item \textbf{The $\Zmod5$ worked case} (\S\ref{sec:z5}, new in v2). The first instance of
        the resulting open problem: the forced floor sharpens to $\muS$, the degree-four sector
        closes at $\gold^{4}$, and the residual contracts to the pentagon-restricted Lehmer
        core (Thms.~\ref{thm:z5floor}--\ref{thm:z5pure}, Prop.~\ref{prop:z5residual}).
  \item \textbf{A commutator scope boundary} (\S\ref{sec:scope}, new in v3). By Shoda's theorem
        the matrix commutator realizes every trace-zero integer matrix, including one whose
        spectrum carries Lehmer's number; hence the floor, gap, and charge structure are
        properties of the abelian semiring only, and no single theorem covers every free
        commutator (Prop.~\ref{prop:commutator}).
\end{itemize}

\subsection{A note on rigor}
Every numerical assertion below was verified in exact arithmetic (\S\ref{sec:method}); no
floating-point comparison was permitted to decide a structural question. We tag each result
with its epistemic status: \forced{} (proved), \computed{} (verified exhaustively over a
stated finite range), or \plausible{} (strongly evidenced, not proved in full generality).
We do \emph{not} claim to resolve Lehmer's problem or the no-Salem conjecture; we
characterize their interaction with the charge character on this substrate.

\section{Preliminaries}\label{sec:prelim}

\subsection{The two operators on the characters}
The operations act on the characters as follows; this table is the algebraic backbone of
everything that follows.

\begin{center}
\begin{tabular}{lll}
\toprule
operation & Character I ($\Mah$) & Character II ($\chi$, in $\Zmod n$) \\
\midrule
$A\dsum B$ & $\Mah(A)\,\Mah(B)$ & multiset union of charges \\
$A\tens B$ & $\prod_{i,j}\max(1,|\alpha_i\beta_j|)$ \ \ (tropical $(\max,+)$) & sumset of charges \\
$\Adams k(A)$ & $\Mah(A)^{k}$ & $\chi\mapsto k\chi \bmod n$ \\
\bottomrule
\end{tabular}
\end{center}

\noindent The $\dsum$ and $\Adams k$ rules for $\Mah$ are immediate from multiplicativity of
$\Mah$ over roots and from $\Adams k$ raising each modulus to the $k$-th power. The $\tens$
rule for $\Mah$ is Proposition~\ref{prop:tensM}.

\subsection{Character I: the Mahler measure and the floor}
\begin{definition}
For monic $p\in\Z[x]$ with roots $\alpha_i$, $\Mah(p)=\prod_i\max(1,|\alpha_i|)$.
\end{definition}

\begin{proposition}[Kronecker]\label{prop:kron}
$\Mah(p)=1$ for a monic $p\in\Z[x]$ if and only if every root of $p$ is a root of unity.
\end{proposition}

We take as given the algebra's defining magnitude property, the \emph{emission gap}.

\begin{proposition}[Emission gap; floor $\gold$]\label{prop:gap}
For charge-admissible objects, $\Mah\in\{1\}\cup[\gold,\infty)$.
\end{proposition}

\noindent\emph{Status (corrected in v2).} \forced{} for the integer-quadratic spectrum (the
gap $(1,\gold)$ is empty among quadratic units, by direct discriminant analysis). In general
the bound $\gold$ is only \plausible{}: what is \emph{forced} is the weaker statement
$\Mah\ge\muS=1.3247\ldots$ on the non-reciprocal class (Lemma~\ref{lem:smyth}) together with
$\Mah\ge\gold^{2}$ on the real part of the reciprocal class (Lemma~\ref{lem:recipunit}). The
band $[\muS,\gold)$ is therefore \emph{not} proven empty; it is part of the residual of
\S\ref{sec:scope}. Lemma~\ref{lem:salem} below shows the sub-$\gold$ algebraic integers that do
exist (Lehmer/Salem numbers) are precisely the charge-inadmissible ones, which is the
structural reason the gap is conjectured to survive the charge restriction.

\begin{definition}[Reciprocal]
$p$ is \emph{reciprocal} if its root multiset is closed under $\alpha\mapsto\alpha^{-1}$,
equivalently $x^{\deg p}p(1/x)=p(x)$. Salem and Lehmer polynomials are reciprocal; the entire
difficulty of Lehmer's problem lives on the reciprocal locus, because Smyth disposes of the
other side outright.
\end{definition}

\begin{lemma}[Smyth]\label{lem:smyth}
Every non-reciprocal monic $p\in\Z[x]$ with $\Mah(p)>1$ satisfies $\Mah(p)\ge\muS=1.3247\ldots$,
the real root of $x^{3}-x-1$ (the plastic number, the smallest Pisot number).
\end{lemma}
\noindent\emph{Status.} \forced{} (Smyth 1971). The bound is unconditional only off the
reciprocal locus---which is exactly where it leaves the residual of \S\ref{sec:z5}.

\begin{lemma}[Real reciprocal units]\label{lem:recipunit}
A reciprocal pair $\{r,r^{-1}\}$ of real conjugates with $r>1$ and rational trace has integer
trace $r+r^{-1}\ge3$ (the value $2$ forces $r=1$), hence $r\ge\gold^{2}=\tfrac{3+\sqrt5}{2}$ and
Mahler contribution $\ge\gold^{2}\approx2.618$.
\end{lemma}
\noindent\emph{Status.} \forced{}.

\subsection{Character II: the charge}
\begin{definition}\label{def:charge}
An object is \emph{charge-admissible} if every conjugate $\alpha$ has $\alpha^{n}\in\R_{>0}$
for some $n\ge1$. The \emph{charge group} is $\Zmod n$ for the least such $n$; equivalently,
writing each conjugate as $|\alpha|e^{2\pi i t}$ with $t\in\Q$, $n$ is the least common
denominator of the $t$'s. The \emph{charge} of $\alpha$ is
$\chi_n(\alpha)=\operatorname{round}(nt)\bmod n\in\Zmod n$.
\end{definition}

\begin{remark}
$\alpha^{n}\in\R_{>0}$ says precisely that $\Arg\alpha\in(2\pi/n)\Z$. The readout
$\chi_n$ is the algebra's phase character; $\Adams n$ \emph{realifies} an object of charge
group $\Zmod n$ (sends every eigenvalue to the positive reals), which is the lever used
repeatedly below.
\end{remark}

\subsection{The golden seed and the \texorpdfstring{$\sqrt5$}{sqrt5} coupling}
The golden ratio $\gold$ is the larger root of $x^2-x-1$; its conjugate is
$\gold'=\tfrac{1-\sqrt5}{2}=-1/\gold\approx-0.618$, satisfying
$\sqrt5=\gold+\gold^{-1}=\gold-\gold'$. The single arithmetic fact that
\[
\boxed{\;\gold'<0,\quad\text{i.e.\ }\gold'\ \text{lies at argument }\pi\;}
\]
drives the parity phenomenon of \S\ref{sec:floor}: the floor-attaining construction must
place $\gold'$ on the charge lattice, and the ray $\pi$ belongs to $(2\pi/n)\Z$ iff $n$ is
even.

\section{The realizable charge groups}\label{sec:structure}

\begin{theorem}[Cyclicity]\label{thm:cyclic}
Every charge group is cyclic. In particular no charge-admissible object has charge group
$\Zmod p\times\Zmod p$, or any non-cyclic abelian group.
\end{theorem}
\begin{proof}
By Definition~\ref{def:charge} the charges of the conjugates are elements
$t\in\Q/\Z\cong\bigoplus_{p}\Z[1/p]/\Z$. They generate a finite subgroup
$G\le\Q/\Z$. Every finite subgroup of $\Q/\Z$ is cyclic: it is the group $\mu_n$ of order
$n=|G|$, since $\Q/\Z$ has exactly one subgroup of each finite order and that subgroup is
$\langle 1/n\rangle\cong\Zmod n$. The charge group is $G$, hence cyclic.
\end{proof}

\noindent\emph{Status.} \forced{}. The content is that the phase character is a rank-one
invariant: each prime contributes at most one independent cyclic label, so the number of
independent primary labels of $\Zmod n$ is $\omega(n)$, the number of distinct prime
divisors of $n$.

\begin{theorem}[Realizability]\label{thm:realize}
For every $n\ge1$, the polynomial $x^{n}-2$ is charge-admissible with charge group $\Zmod n$,
attains all $n$ charges, and has $\Mah=2$. Hence the realizable charge groups are exactly
$\{\Zmod n: n\ge1\}$.
\end{theorem}
\begin{proof}
The roots are $2^{1/n}\zeta^{k}$, $\zeta=e^{2\pi i/n}$, $k=0,\dots,n-1$, at arguments
$2\pi k/n$ exactly filling $(2\pi/n)\Z/2\pi\Z$; the least common denominator is $n$, so the
charge group is $\Zmod n$ and every charge $k$ occurs. Each root has modulus $2^{1/n}>1$, so
$\Mah=\prod_k 2^{1/n}=2$. Combined with Theorem~\ref{thm:cyclic}, the realizable groups are
the finite cyclic groups.
\end{proof}
\emph{Status.} \forced{}; verified for $n=3,\dots,7$ (App.~\ref{app:ledger}, entry A).

\begin{theorem}[Primary decomposition via Adams operations]\label{thm:crt}
Let $n=\prod_{i} p_i^{e_i}$. Then $\Zmod n\cong\prod_i\Zmod{p_i^{e_i}}$ by the Chinese
Remainder Theorem, and this decomposition is realized inside the algebra: the Adams operation
$\Adams{\,n/p_i^{e_i}}$ annihilates the prime-to-$p_i$ part of the charge and projects onto
the $p_i$-primary component $\Zmod{p_i^{e_i}}$.
\end{theorem}
\begin{proof}
On charges $\Adams k$ acts as multiplication by $k$ on $\Zmod n$
(table, \S\ref{sec:prelim}). For $k=n/p_i^{e_i}$, multiplication by $k$ has image the unique
subgroup of order $p_i^{e_i}$ and kernel the complementary factor, i.e.\ it is the CRT
projection onto $\Zmod{p_i^{e_i}}$. Realification: applying $\Adams{k}$ to an object of charge
group $\Zmod n$ yields one whose residual charge group is exactly $\Zmod{p_i^{e_i}}$.
\end{proof}
\emph{Status.} \forced{} (CRT). Verified on $x^6-2$: $\Adams3$ leaves residual group
$\Zmod2$, $\Adams2$ leaves $\Zmod3$ (App.~\ref{app:ledger}, entry C). Thus prime powers give a
\emph{nested} filtration (via $\Adams{p^{j}}$) and distinct primes a \emph{flat} product,
with $\omega(n)$ independent labels.

\section{Composition under the tensor product}\label{sec:compose}

\begin{proposition}[Measure under $\tens$; tropical law]\label{prop:tensM}
With $A,B$ having roots $\{\alpha_i\},\{\beta_j\}$,
\[
\Mah(A\tens B)=\prod_{i,j}\max\bigl(1,|\alpha_i\beta_j|\bigr),
\qquad\text{equivalently}\qquad
\log\Mah(A\tens B)=\sum_{i,j}\bigl(\log|\alpha_i|+\log|\beta_j|\bigr)^{+},
\]
where $t^{+}=\max(0,t)$. This is the tropical $(\max,+)$ law on the log-moduli. It does
\emph{not} factor as $\Mah(A)^{\deg B}\Mah(B)^{\deg A}$ in general; the two agree exactly when
no product $\alpha_i\beta_j$ crosses the unit circle (e.g.\ when every $|\alpha_i|,|\beta_j|>1$),
and disagree otherwise.
\end{proposition}
\begin{proof}
The eigenvalues of $A\tens B$ are the $\deg A\cdot\deg B$ products $\alpha_i\beta_j$, so by the
definition of $\Mah$, $\Mah(A\tens B)=\prod_{i,j}\max(1,|\alpha_i\beta_j|)$; taking logarithms
gives the stated sum. The factored form would require
$\max(1,|\alpha_i\beta_j|)=\max(1,|\alpha_i|)\max(1,|\beta_j|)$ for every pair, which fails
precisely when one factor lies inside the unit circle and the other outside with the product on
the opposite side: e.g.\ $\gold\tens\gold$ has eigenvalues $\{\gold^{2},-1,-1,\gold'^{2}\}$,
giving $\Mah=\gold^{2}$, whereas $\Mah(\gold)^{2}\Mah(\gold)^{2}=\gold^{4}$
(App.~\ref{app:ledger}, entry M).
\end{proof}

\begin{remark}
This corrects an error in earlier drafts (v1--v3), which stated the factored form. The
correct tropical law is the one carried by the corpus's operator-algebra specification (S4);
no result of this paper depends on the tensor measure---the $\operatorname{lcm}$ law
(Thm.~\ref{thm:lcm}) and the parity and $\Zmod5$ proofs use the \emph{charge} character and
the constructions $q_k$, $x^{n}-2$, and direct quartics---so the correction is confined to
this proposition and the backbone table.
\end{remark}

\begin{theorem}[The $\operatorname{lcm}$ law]\label{thm:lcm}
If $A,B$ have charge groups $\Zmod m,\Zmod k$, then $A\tens B$ has charge group
$\Zmod{\operatorname{lcm}(m,k)}$. If $\gcd(m,k)=1$ this is the full product
$\Zmod m\times\Zmod k\cong\Zmod{mk}$ and the pair $(\chi_m,\chi_k)$ is recoverable from the
composite charge; if $\gcd(m,k)>1$, the $\gcd$-part collapses.
\end{theorem}
\begin{proof}
Arguments add under multiplication of eigenvalues: a root of $A$ at $2\pi a/m$ and a root of
$B$ at $2\pi b/k$ produce a root of $A\tens B$ at $2\pi(a/m+b/k)$. The set
$\{a/m+b/k\bmod 1\}$ generates the subgroup $\tfrac1m\Z+\tfrac1k\Z=\tfrac1{\operatorname{lcm}(m,k)}\Z$
of $\Q/\Z$, whose order is $\operatorname{lcm}(m,k)$. Coprimality is the CRT case
$\Zmod{mk}\cong\Zmod m\times\Zmod k$, in which the projections recover $(\chi_m,\chi_k)$.
\end{proof}

\begin{theorem}[Safe composition $\Leftrightarrow$ coprime]\label{thm:safe}
Two charge sectors compose under $\tens$ without making any previously forbidden charge
reachable if and only if their orders are coprime.
\end{theorem}
\begin{proof}
By Theorem~\ref{thm:lcm}, coprime orders yield $\Zmod m\times\Zmod k$, on which the two
charges remain independent coordinates; no label of one sector is produced by the other, so
the forbidden sets stay separated (lossless). If $\gcd(m,k)=d>1$, the composite group has
order $\operatorname{lcm}(m,k)=mk/d<mk$; the surjection
$\Zmod m\times\Zmod k\twoheadrightarrow\Zmod{\operatorname{lcm}}$ has nontrivial kernel, and
its identifications can map an allowed pair onto a forbidden composite charge---a leak.
\end{proof}
\emph{Status.} \forced{}. Illustration: $x^3-2\ \tens\ x^4-2$ has $\gcd(3,4)=1$; the engine
returns charpoly $x^{12}-128$, charge group $\Zmod{12}=\Zmod{\operatorname{lcm}(3,4)}$, and
$\Mah=128=2^{7}$ (all twelve products have modulus $2^{7/12}>1$, so the tropical law of
Prop.~\ref{prop:tensM} reduces to $(2^{7/12})^{12}=2^{7}$; this is an off-circle case where the
factored form happens to agree) (App.~\ref{app:ledger}, entry D).

\section{The coupling: a parity-graded floor}\label{sec:floor}

We now address the coupling question. For $n\ge1$ define
\[
\mu(n)\;=\;\inf\bigl\{\,\Mah(O) : O\ \text{charge-admissible},\ \text{charge group }\Zmod n,\ \Mah(O)>1\,\bigr\}.
\]

\begin{theorem}[Parity-graded floor]\label{thm:parity}
\[
\mu(n)=
\begin{cases}
\gold, & n\ \text{even},\\[2pt]
2, & n\ \text{odd},
\end{cases}
\]
Here $\mu(n)=2$ for odd $n$ is the \emph{realized} floor; the \emph{forced} lower bound for odd
$n$ is only $\muS$ (Lemma~\ref{lem:smyth}), and the equality $\mu(\mathrm{odd})=2$ is held at
\computed{} strength (\S\ref{sec:z5}). The bound $\gold$ is attained \emph{iff} $n$ is even.
Symmetrically, $\mu(\mathrm{even})=\gold$ is two statements of unequal strength:
$\mu(\mathrm{even})\le\gold$ is \forced{} (the construction below attains it), while
$\mu(\mathrm{even})\ge\gold$ inherits the \plausible{} status of the universal bound
(Prop.~\ref{prop:gap}). The even floor is attained constructively: for $n=2k$,
\[
q_k(x)\;=\;x^{2k}+x^{k}-1
\]
is charge-admissible with charge group $\Zmod{2k}$, attains all $2k$ charges, and has
$\Mah(q_k)=\gold$.
\end{theorem}

\begin{proof}[Proof of even attainment]
Substitute $y=x^{k}$ in $q_k$: $y^2+y-1=0$ gives $y=\gold-1=1/\gold>0$ and
$y=-\gold<0$.
\begin{itemize}[nosep]
  \item $x^{k}=1/\gold$ contributes $k$ roots of modulus $(1/\gold)^{1/k}<1$ at arguments
        $2\pi j/k$ (even multiples of $\pi/k$); these lie inside the unit circle and add
        nothing to $\Mah$.
  \item $x^{k}=-\gold$ contributes $k$ roots of modulus $\gold^{1/k}>1$ at arguments
        $(2j+1)\pi/k$ (odd multiples of $\pi/k$); their product of moduli is
        $(\gold^{1/k})^{k}=\gold$.
\end{itemize}
Hence $\Mah(q_k)=\gold$. The full set of arguments is all multiples of $\pi/k$, i.e.\
$(2\pi/2k)\Z$, with least common denominator $2k$; so the charge group is $\Zmod{2k}$ and all
charges occur. For $k=2$ this is the factorization
$x^4+x^2-1=(x^2+\gold)(x^2+\gold')$, placing the measure-bearing pair $\pm i\sqrt{\gold}$ on
the imaginary axis and the inert pair $\pm\sqrt{1/\gold}$ on the real axis.
\end{proof}

\begin{lemma}[The $\pi$-ray obstruction]\label{lem:pi}
The golden conjugate $\gold'=\tfrac{1-\sqrt5}{2}<0$ lies at argument $\pi$. The ray $\pi$
belongs to the lattice $(2\pi/n)\Z$ if and only if $n$ is even. Consequently the golden
seed---the unique structure realizing the universal bound $\Mah=\gold$ at nontrivial
charge---can be hosted on a charge lattice $\Zmod n$ only when $n$ is even.
\end{lemma}
\begin{proof}
$\pi=(2\pi/n)\cdot m$ for an integer $m$ iff $n=2m$, i.e.\ $n$ even. Any object attaining
$\Mah=\gold$ through the golden structure carries a conjugate at argument $\pi$ (the image of
$\gold'$ under the construction); if its charge group were $\Zmod n$ with $n$ odd, that
conjugate would force $\pi\in(2\pi/n)\Z$, a contradiction.
\end{proof}

\begin{proof}[Proof of the odd floor (forced part) and status of the rest]
By Theorem~\ref{thm:realize}, $x^{n}-2$ realizes $\Zmod n$ with $\Mah=2$, so $\mu(n)\le2$. For
the lower bound, what is \emph{forced} is $\mu(n)\ge\muS$: a non-reciprocal object obeys
Lemma~\ref{lem:smyth}, and a reciprocal one obeys $\Mah\ge\gold^{2}>2$ on its real part
(Lemma~\ref{lem:recipunit})---the two cases worked at $\Zmod5$ in \S\ref{sec:z5}.
Lemma~\ref{lem:pi} blocks the golden construction for odd $n$, so $\gold$ is unattained; and
Theorem~\ref{thm:salemgap} gives the \computed{} statement that no admissible measure lies in
$(\gold,2)$. The equality $\mu(\mathrm{odd})=2$ thus rests on emptiness of the non-reciprocal
class in $[\muS,2)$, verified over the window but not proved (\S\ref{sec:z5},\,\S\ref{sec:scope}).
The even case combines the construction above with the universal bound $\gold$.
\end{proof}

\noindent\emph{Status.} \forced{} for the even attainment (constructive), the $\pi$-ray
obstruction, and the forced odd lower bound $\muS$. The odd \emph{realized} floor
$\mu(\mathrm{odd})=2$ is \forced{} for $\Zmod3$ (Thm.~\ref{thm:salemgap}) and \computed{} in
general. Verified: $q_k$ for $k=2,\dots,5$ realizes $\Zmod{2k}$ at $\Mah=\gold$
(App.~\ref{app:ledger}, entry E).

\begin{theorem}[Totally-positive gap; the blocked cheap route]\label{thm:tp}
A totally positive algebraic integer (all conjugates real and positive) has no Mahler measure
in $(1,2)$; the least value exceeding $1$ is exactly $2$, attained by $(x-1)(x-2)$.
Consequently the cheapest charging route---$\zeta_n$ (a primitive $n$-th root of unity,
$\Mah=1$) tensored with a positive-real measure-bearing factor---cannot produce an odd-charge
object with $\Mah\in(\gold,2)$.
\end{theorem}
\emph{Status.} \computed{} (verified for all such integers of degree $\le7$). The relevance:
to carry an odd charge cheaply one multiplies the free cyclotomic charge $\zeta_n$
($\Mah=1$) by a factor whose roots avoid the argument $\pi$ (else the lattice doubles to even
order); such a factor is totally positive, and the gap forbids its measure from entering
$(\gold,2)$. The next totally positive value after $2$ is $\gold^{2}\approx2.618$
(App.~\ref{app:ledger}, entry G).

\section{The odd-charge gap is the no-Salem closure}\label{sec:salem}

\begin{lemma}[Salem exclusion]\label{lem:salem}
A Salem number carries no finite charge group. More generally, any algebraic integer with an
irrational conjugate angle is charge-inadmissible.
\end{lemma}
\begin{proof}
Let $\tau$ be a Salem number: $\tau>1$, $1/\tau$ its only other real conjugate, and all
remaining conjugates on $|z|=1$. Those unit-circle conjugates are roots of the minimal
polynomial of $\tau$, which is irreducible and non-cyclotomic; hence none of them is a root
of unity (a root of unity has cyclotomic minimal polynomial, which would have to coincide
with the Salem polynomial). A unit-circle algebraic integer that is not a root of unity has
argument an irrational multiple of $2\pi$, so $\tau^{n}\notin\R_{>0}$ for every $n$:
$\tau$ is charge-inadmissible. The general statement is immediate from
Definition~\ref{def:charge}.
\end{proof}

\begin{theorem}[Odd gap $=$ no-Salem closure]\label{thm:salemgap}
For odd $n$, the open interval $(\gold,2)$ contains no Mahler measure of any
charge-admissible object with charge group $\Zmod n$. Moreover, the algebraic integers with
$\Mah\in(\gold,2)$ lying near odd-order lattices are exactly Salem and Pisot numbers with
irrational angles; by Lemma~\ref{lem:salem} these are charge-inadmissible. Thus the
odd-charge gap $(\gold,2)$ coincides with the no-Salem closure of the algebra: the measures
that would fill it are precisely the excluded Salem occupants.
\end{theorem}

\begin{proof}[Proof for $\Zmod3$ (the forced case)]
An irreducible cubic on the $\Zmod3$ lattice has signature $(1,1)$: a real root at argument
$0$ and a conjugate pair at $\pm120^\circ$, with quadratic factor
$x^2-2\rho\cos120^\circ\,x+\rho^2=x^2+\rho x+\rho^2$ since $\cos120^\circ=-\tfrac12$.
Integrality of the coefficients forces $\rho,\rho^2\in\Z$, hence $\rho\in\Z$. Writing the
cubic as $(x-\alpha)(x^2+\rho x+\rho^2)$ with $\alpha\in\R$, the coefficient relations give
$c_1=\rho(\rho-\alpha)$ and $c_2=\rho-\alpha$, so $c_1=\rho\,c_2$ with $\rho\in\Z$; the only
sub-$2$ solution is $x^3-2$ ($\rho=0$, $\Mah=2$). No configuration yields
$\Mah\in(1,2)$; in particular $(\gold,2)$ is empty for $\Zmod3$. Exhaustive search over
$|c_i|\le6$ confirms zero exceptions.
\end{proof}

\noindent\emph{General evidence and status.} A two-stage search---a fast double-precision
screen followed by exact, high-precision re-classification with tolerance matched to the
working precision---was run over monic integer polynomials (degree $\le7$, small
coefficients), reciprocal/Salem-adjacent polynomials (degree $\le12$), and the
totally-positive family (degree $\le7$). It returned \emph{no} genuine odd-charge object with
$\Mah\in(\gold,2)$. Every one of the $17$ apparent sub-$2$ candidates surfaced by the loose
screen had, on exact re-classification, \emph{no} finite charge group: each is an
irrational-angle Salem or Pisot number. The keystone is
\[
\beta_4=x^4-x^3-x^2-x+1,\qquad \Mah(\beta_4)=1.72208\ldots\in(\gold,2),
\]
the minimal degree-$4$ Salem number, correctly classified as charge-inadmissible
(App.~\ref{app:ledger}, entry F). Thus the would-be occupants of the odd-charge gap are the
exact objects excluded by the emission gap. The statement is \forced{} for $\Zmod3$ and via
Lemma~\ref{lem:salem}, and \computed{}/\plausible{} for the all-degrees claim, which would
require a Mahler coefficient-bound argument rather than a finite search.

\begin{corollary}
The parity dichotomy of Theorem~\ref{thm:parity} and the Salem gap of
Theorem~\ref{thm:salemgap} are two faces of one fact: the universal magnitude floor $\gold$
is realizable at nontrivial charge \emph{iff} the charge lattice contains the argument $\pi$
(i.e.\ $n$ even); when it does not (odd $n$), the next admissible magnitude is $2$, and the
band $(\gold,2)$ is occupied only by the charge-inadmissible Salem spectrum.
\end{corollary}

\section{The \texorpdfstring{$\Zmod5$}{Z/5} case: first worked instance of the open problem}\label{sec:z5}

The open problem of \S\ref{sec:scope}---prove for all degrees that odd charge admits no measure
in $(\gold,2)$---is here worked for the smallest odd modulus on which the elementary $\Zmod3$
argument fails. The $\Zmod3$ proof closes because $2\cos120^\circ=-1\in\Z$; at $\Zmod5$ the
relevant cosines are irrational and Galois-coupled, and the case splits into a \emph{forced}
skeleton and a single \emph{computed} residual which is exactly Lehmer's problem restricted to
the pentagon lattice. Nothing here closes that residual; the contribution is to sharpen the
forced floor, close the degree-four sector, and name what remains.

\begin{proposition}[Pentagon cosines]\label{prop:pentcos}
$2\cos72^\circ=\gold-1$ and $2\cos144^\circ=-\gold$ are the two (irrational) roots of
$x^2+x-1$, hence Galois-conjugate over $\Q$ under $\sqrt5\mapsto-\sqrt5$.
\end{proposition}
\noindent\emph{Status.} \forced{} (symbolic). A single $\pm72^\circ$ pair of rational modulus
$s$ has factor $x^2-s(\gold-1)x+s^2$ with \emph{irrational} linear term; integrality is
recoverable only by adjoining the $\pm144^\circ$ Galois partner so the $\sqrt5$-parts cancel.
This forced coupling is the feature $\Zmod3$ never needs.

\begin{theorem}[Forced floor at $\Zmod5$]\label{thm:z5floor}
Let $O$ be charge-admissible with charge group $\Zmod5$ and $\Mah(O)>1$.
\begin{enumerate}[label=\emph{(\alph*)},nosep]
  \item If $O$ is non-reciprocal, then $\Mah(O)\ge\muS=1.3247\ldots$ \forced{}
  \item If $O$ is reciprocal, then $O=\Phi_5^{\,a}\cdot U$ with $U$ a product of reciprocal
        units off the unit circle; its real part obeys $\Mah\ge\gold^{2}$, and no reciprocal
        charge-$\Zmod5$ object in the verification window has $\Mah<\gold^{2}$.
        \forced{} (real part) / \computed{} (window)
\end{enumerate}
This improves the realification bound $2^{1/5}=1.1487$---from $\Adams5O$ totally positive with
$\Mah(\Adams5O)=\Mah(O)^{5}\ge2$ (Thm.~\ref{thm:tp})---to $\muS$.
\end{theorem}
\begin{proof}
(a) is Lemma~\ref{lem:smyth}. (b) By Lemma~\ref{lem:salem} an on-unit root of a
charge-admissible non-cyclotomic object is a fifth root of unity; strip the resulting
$\Phi_5^{\,a}$. The remainder $U$ is reciprocal, charge-admissible, hence not Salem, so it has
no on-circle conjugate at all; every root lies off the circle paired with its inverse. The real
such pairs are real reciprocal units, $\Mah\ge\gold^{2}$ by Lemma~\ref{lem:recipunit}; the
complex such pairs sit at $\pm72^\circ,\pm144^\circ$, coupled in $\Q(\sqrt5)$ as in
Theorem~\ref{thm:pentagon}. The window (App.~\ref{app:ledger}, entries I--K) exhibits only
$\Phi_5^{\,a}\times(\text{real reciprocal units})$ with $\Mah\in\{1,\gold^{2},2+\sqrt3,\dots\}$,
all $\ge\gold^{2}>2$.
\end{proof}

\begin{theorem}[Pure pentagon: the degree-four sector]\label{thm:pentagon}
Every irreducible quartic with charge group $\Zmod5$ has $\Mah\in\{1\}\cup[\gold^{4},\infty)$,
$\gold^{4}\approx6.854$; in particular $\Mah\notin(1,\gold^{4})$. \forced{}
\end{theorem}
\begin{proof}
By Lemma~\ref{lem:salem} the object has no real root ($180^\circ\notin$ the lattice), so it is a
$\pm72^\circ$ pair of modulus $s$ times a $\pm144^\circ$ pair of modulus $t$:
$O=\bigl(x^2-s(\gold-1)x+s^2\bigr)\bigl(x^2+t\gold\,x+t^2\bigr)$. Integrality of all four
coefficients forces the $\sqrt5$-parts to cancel, i.e.\ $t=\sigma(s)$ under
$\sqrt5\mapsto-\sqrt5$; equivalently $s,t$ are the two positive roots of a totally real integer
quadratic $y^2-ky+m$, $k=s+t\in\Z$, $m=st\in\Z_{>0}$. With two roots at $s$ and two at $t$,
$\Mah=\max(1,s)^2\max(1,t)^2$. The case analysis on $(k,m)$ bottoms out at $(k,m)=(3,1)$:
$s=\gold^{2}$ at $\pm72^\circ$, $t=\gold^{-2}$ at $\pm144^\circ$, $\Mah=\gold^{4}$, with explicit
minimizer
\[
x^4-x^3+6x^2+4x+1,\qquad \text{charge}=\Zmod5,\quad \Mah=\gold^{4}.
\]
Irreducibility forces this Galois-coupled form for \emph{every} charge-$\Zmod5$ quartic (a
$72^\circ$ root drags in its $\sqrt5$-partner at $144^\circ$), so the classification is
complete; an exhaustive scan ($|c|\le10$) returns distinct measures exactly $\{1,\gold^{4}\}$
(App.~\ref{app:ledger}, entry J).
\end{proof}
\begin{remark}
The contrast with $\Zmod3$ is now exact: $\Zmod3$ couples nothing and bottoms out at $2$;
$\Zmod5$ couples the two moduli as $\Q(\sqrt5)$-conjugates and bottoms out at $\gold^{4}$. The
coupling that makes the odd case hard \emph{in general} is the very thing that widens the
degree-four gap here.
\end{remark}

\begin{theorem}[Pure power realizes the floor]\label{thm:z5pure}
For $m\ge2$, $x^{5}-m$ has charge group $\Zmod5$, is non-reciprocal, and has $\Mah=m$. Hence the
family realizes $\Mah\in\{1\}\cup[2,\infty)$, attaining $\mu(5)=2$ at $x^{5}-2$. \forced{}
\end{theorem}
\begin{proof}
Roots $m^{1/5}\zeta^{k}$, all of modulus $m^{1/5}>1$, give charge $\Zmod5$ and $\Mah=m$. Closure
under $\alpha\mapsto\alpha^{-1}$ would require $m^{-1/5}$ among the roots, which it is not; hence
non-reciprocal (consistent with Theorem~\ref{thm:z5floor}(a), $m\ge2>\muS$). The real root
$m^{1/5}$ is Galois-conjugate to the four complex roots and cannot be split off, so $x^{5}-m$ is
genuinely quintic and escapes the degree-four analysis of Theorem~\ref{thm:pentagon}.
\end{proof}

\begin{proposition}[The residual, named]\label{prop:z5residual}
Combining Theorems~\ref{thm:z5floor}--\ref{thm:z5pure}: the reciprocal charge-$\Zmod5$ class
lies in $\{1\}\cup[\gold^{2},\infty)$ (above $2$); the non-reciprocal class is forced into
$\{1\}\cup[\muS,\infty)$ with realized minimum $2$. Hence a charge-$\Zmod5$ object can land in
$(1,2)$ \emph{only} by being non-reciprocal with $\Mah\in[\muS,2)$. Over the verification window
(quartics $|c|\le10$, quintics $|c|\le4$, sextics $|c|\le3$) this set is empty, the boundary
object $x^{5}-2$ sitting at exactly $2$. \computed{}
\end{proposition}
\begin{remark}[Why this does not promote]
Three handles reach the residual and none reaches $2$. \emph{Realification}: $\Adams5O$ is
totally positive with $\Mah(\Adams5O)=\Mah(O)^{5}$, giving only $\Mah(O)\ge2^{1/5}$; asking
$\Adams5O$ to clear $32$ rather than $2$ is the original claim restated, hence circular.
\emph{Smyth}: Lemma~\ref{lem:smyth} forces $\muS$, the best off-the-shelf bound, still below
$2$. \emph{Classification}: Theorems~\ref{thm:pentagon} and~\ref{thm:z5pure} close degree four
and the pure-power family, but a general non-reciprocal charge-$\Zmod5$ object of degree $\ge5$
with non-constant modulus structure is unclassified. The residual---``no non-reciprocal
charge-$\Zmod5$ object has $\Mah\in[\muS,2)$''---is the reciprocal/Salem core of Lehmer's
problem specialized to the pentagon lattice, equivalently the totally-positive $(1,2)$ gap
surviving fifth-root pullback with the rotation-class coupling that $\Adams5$ destroys. That is the
Schur--Siegel--Smyth (SSS) trace problem (Siegel~1945; Smyth~1984), a recognized open problem: the
residual \emph{reduces into} SSS, and whether it is \emph{equivalent} to SSS or only \emph{embeds
into} it---the exact relationship (one-way reduction vs.\ equivalence)---is itself an explicit open
sub-question, not closable by the methods here. By the
asymmetry rule (a finite search may extend a claim to \computed{} but never promote it to
\forced{}), the four forced sub-results plus a clean window are the honest maximum:
\textbf{the open problem of \S\ref{sec:scope} is refined at $\Zmod5$, not closed.} That the
residual resists matrix-algebra structure is no accident: by Proposition~\ref{prop:commutator}
the Lehmer/Salem core it names is realizable outright as a commutator spectrum, so only the
charge restriction---not the ambient algebra---can do the excluding.
\end{remark}

\section{Methodology}\label{sec:method}

\paragraph{Exact arithmetic.} Charge-group membership is the assertion
$\alpha^{n}\in\R_{>0}$, a decision about whether an imaginary part is exactly zero. We never
let machine doubles adjudicate it. Roots were computed at $25$--$35$ decimal digits with
$\mathtt{mpmath}$, multiplicities removed via the squarefree part, and the membership
tolerance set well below the precision floor but far above true-zero residuals
($10^{-18}$ at $30$ digits). The one recurring failure mode in this line of work was a loose
angular tolerance on double precision mis-binning an irrational angle as on-lattice; the
discipline above is its antidote, and it is what distinguishes the spurious screen
candidates of \S\ref{sec:salem} from genuine witnesses.

\paragraph{Engine validation.} The operations were checked against an independent
integer-matrix engine in which $\dsum$ is the matrix direct sum, $\tens$ the Kronecker
product, and $\Adams k$ the matrix $k$-th power; charge readouts were taken from companion
matrices of the relevant polynomials. The consolidation ledger of
Appendix~\ref{app:ledger} was produced by this engine and underlies every numerical claim in
the text.

\paragraph{Epistemic tagging.} Each result carries one of \forced{}/\computed{}/\plausible{}.
This is not decoration: the parity floor's even half is a constructive proof, its odd half is
proved only for $\Zmod3$ and verified to a finite degree beyond, and the universal lower
bound itself is proved only for quadratics. Conflating these would misrepresent what the
algebra is known to do.

\section{Scope, limitations, and open problems}\label{sec:scope}

\begin{center}
\begin{tabular}{lll}
\toprule
Result & Statement & Status \\
\midrule
Cyclicity (Thm.~\ref{thm:cyclic}) & charge groups are cyclic & \forced{} \\
Realizability (Thm.~\ref{thm:realize}) & every $\Zmod n$ occurs & \forced{} \\
CRT / Adams (Thm.~\ref{thm:crt}) & primary projectors $\Adams{n/p^{e}}$ & \forced{} \\
$\operatorname{lcm}$ law (Thm.~\ref{thm:lcm}) & $\tens$ gives $\Zmod{\operatorname{lcm}}$ & \forced{} \\
Safe $\Leftrightarrow$ coprime (Thm.~\ref{thm:safe}) & lossless iff coprime orders & \forced{} \\
Even floor, attainment (Thm.~\ref{thm:parity}) & $\mu(\text{even})\le\gold$ ($q_k$ attains) & \forced{} \\
Even floor, lower bound (Thm.~\ref{thm:parity}) & $\mu(\text{even})\ge\gold$ & \plausible{} (Prop.~\ref{prop:gap}) \\
Forced odd floor (Lem.~\ref{lem:smyth}) & $\Mah\ge\muS$, odd non-reciprocal & \forced{} (Smyth) \\
Odd realized floor (Thm.~\ref{thm:parity}) & $\mu(\text{odd})=2$ & \forced{} ($\Zmod3$) / \computed{} \\
$\Zmod5$ degree-4 (Thm.~\ref{thm:pentagon}) & $\Mah\in\{1\}\cup[\gold^{4},\infty)$ & \forced{} \\
$\Zmod5$ pure power (Thm.~\ref{thm:z5pure}) & $x^{5}-m$: $\Mah=m$; $\mu(5)=2$ at $x^{5}-2$ & \forced{} \\
$\Zmod5$ residual (Prop.~\ref{prop:z5residual}) & non-recip in $[\muS,2)$ empty over window & \computed{} \\
Totally-positive gap (Thm.~\ref{thm:tp}) & no measure in $(1,2)$ & \computed{} (deg $\le7$) \\
Salem gap (Thm.~\ref{thm:salemgap}) & $(\gold,2)$ empty for odd $n$ & \forced{} ($\Zmod3$) / \plausible{} \\
Universal bound (Prop.~\ref{prop:gap}) & forced floor $\muS$; $\gold$ conjectural & \forced{} ($\muS$) / \plausible{} ($\gold$) \\
Commutator boundary (Prop.~\ref{prop:commutator}) & no uniform theorem over free commutators & \forced{} (Shoda) \\
\bottomrule
\end{tabular}
\end{center}

\subsection*{The commutator boundary}
Every result above governs the abelian spectral semiring---objects built from
$\dsum,\tens,\Adams n$, on which the two characters transform predictably. The natural
non-abelian operation, the matrix commutator $[A,B]=AB-BA$, lies outside it, and is the exact
boundary of the theory.

\begin{proposition}[Commutator escape]\label{prop:commutator}
The matrix commutator is not a semiring operation, and no floor, gap, or charge constraint
extends across it. Concretely, by Albert--Muckenhoupt a matrix over a field is a commutator
$[A,B]$ if and only if it has trace zero, and every trace-zero integer matrix is a commutator
of integer matrices (Laffey--Reams). The companion matrix of $L(x)\,(x-1)$, where
$L$ is Lehmer's polynomial, is a trace-zero integer matrix whose characteristic polynomial
carries Lehmer's number $\tau$. Since $\tau$ is a Salem number, this commutator has
\[
\Mah=\Mah(L)\,\Mah(x-1)=1.17628\ldots\in(1,\gold)
\qquad\text{and}\qquad
\text{charge group }=\ \bot,
\]
i.e.\ it lies \emph{below} the emission floor \emph{and} off every angle lattice
simultaneously. Consequently:
\begin{enumerate}[label=\emph{(\roman*)},nosep]
  \item the emission gap (Prop.~\ref{prop:gap}) and the charge structure of
        \S\ref{sec:structure} do not extend to commutators;
  \item \emph{no single theorem governs every free commutator}: any uniform floor/gap/charge
        statement fails on this witness, so the goal is unsatisfiable;
  \item the no-Salem closure (Thm.~\ref{thm:salemgap}) is \emph{operation-relative}---Salem
        numbers are excluded by the restriction to the charge-admissible abelian semiring, not
        by the ambient matrix algebra, which realizes them freely through the commutator.
\end{enumerate}
\end{proposition}
\begin{proof}
The first two sentences are Albert--Muckenhoupt and Laffey--Reams. Lehmer's polynomial
$L=x^{10}+x^{9}-x^{7}-x^{6}-x^{5}-x^{4}-x^{3}+x+1$ has trace $-1$, so
$L(x)(x-1)=x^{11}-x^{9}-x^{8}+x^{3}+x^{2}-1$ has trace $0$; its companion matrix $C$ is thus a
trace-zero integer matrix and a commutator, with characteristic polynomial $L(x)(x-1)$. Hence
the spectrum of $C$ contains Lehmer's number $\tau$ together with $1$. As $\Mah$ is
multiplicative and $\Mah(x-1)=1$, $\Mah(C)=\Mah(L)=\tau=1.17628\ldots\in(1,\gold)$. By
Lemma~\ref{lem:salem} the on-circle conjugates of $\tau$ are at irrational angles, so $C$ is
charge-inadmissible. The three conclusions follow because $C$ is a commutator violating each
constraint.
\end{proof}

\noindent\emph{Status.} \forced{}. This is the sharp form of the maxim that the commutator is
``the one door'': Shoda fixes the door's width at the full trace-zero set, and Lehmer's number
witnesses that the door admits the entire sub-$\gold$ Salem spectrum. It sharpens the
confinement reading of the programme---the abelian semiring is the box, the commutator is the
door---without altering any result of \S\S\ref{sec:structure}--\ref{sec:z5}.
\emph{all} degrees, that no charge-admissible object with odd charge group has
$\Mah\in(\gold,2)$. A finite search cannot settle this: Mahler's coefficient bound permits
larger coefficients at fixed measure, so a complete proof must proceed structurally---most
plausibly by upgrading Lemma~\ref{lem:salem} to a full dichotomy (every algebraic integer with
$\Mah\in(\gold,2)$ is either Salem-type with irrational angle, or has even charge group). The
$\Zmod5$ worked case (\S\ref{sec:z5}) is the current state of that program: it forces the floor
up to $\muS$, closes the degree-four sector at $\gold^{4}$, and contracts the residual to a
single named set---the non-reciprocal charge-$\Zmod5$ objects in $[\muS,2)$, which is the
reciprocal/Salem core of Lehmer's problem on the pentagon lattice, equivalently the
totally-positive $(1,2)$ gap surviving fifth-root pullback. That residual is trace-problem
territory; the same structural upgrade would promote the universal bound (Prop.~\ref{prop:gap}),
the general odd floor, and $\mu(5)=2$ to \forced{}. Until then the problem is \emph{refined and
still open}.

\paragraph{Caveat on the larger programme.} The substrate's deeper conjectures---a resolution
of Lehmer's problem, or particular continuum couplings posited elsewhere in the
project---are \emph{not} touched here and remain open or conditional. What this note
establishes is internal to the charge--measure interaction on the abelian spectral semiring:
the algebra's two characters are orthogonal as readouts but coupled at the extremal floor, with
the coupling governed by the single parity bit ``does the charge lattice contain $\pi$.'' The
boundary of that statement is Proposition~\ref{prop:commutator}: pass to the non-abelian
commutator and every constraint dissolves, Lehmer's number included. Confinement, here, is
exactly the choice to stay abelian.

\appendix
\section{Consolidation ledger}\label{app:ledger}
The following were regenerated in exact arithmetic immediately prior to writing; each grounds
a claim in the text.

\begin{center}
\begin{tabular}{cll}
\toprule
& object & exact result \\
\midrule
A & $x^{n}-2,\ n=3\text{--}7$ & charge group $\Zmod n$, all charges, $\Mah=2$ \\
B & $x^{4}-2$ & charge group $\Zmod4$ (single generator), not $\Zmod2\times\Zmod2$ \\
C & $x^{6}-2$ & $\Zmod6$; $\Adams3$ residual $\Zmod2$, $\Adams2$ residual $\Zmod3$ \\
D & $x^{3}-2\ \tens\ x^{4}-2$ & charpoly $x^{12}-128$, charge $\Zmod{12}$, $\Mah=128=2^{7}$ \\
E & $x^{2k}+x^{k}-1,\ k=2\text{--}5$ & charge $\Zmod{2k}$, all charges, $\Mah=\gold$ \\
F & $\beta_4=x^4-x^3-x^2-x+1$ & $\Mah=1.72208\in(\gold,2)$, charge group \emph{none} (irrational) \\
G & $(x-1)(x-2)$ & totally positive, $\Mah=2$; next value $\gold^2\approx2.618$ \\
H & --- & $\gold=1.61803$; even floor $\gold$, odd floor $2$, global admissible min $\gold$ \\
\midrule
I & $x^4-x^3+6x^2+4x+1$ & charge $\Zmod5$, $\Mah=\gold^{4}=6.854$, irreducible, non-reciprocal \\
J & charge-$\Zmod5$ quartics, $|c|\le10$ & distinct $\Mah=\{1,\gold^{4}\}$; none in $(1,\gold^{4})$ \\
K & $\Phi_5\!\cdot\!(x^2-3x+1)$;\ $x^{5}-2$ & recip.\ $\Mah=\gold^{2}$;\ \ $x^{5}-2$ non-recip.\ $\Mah=2=\mu(5)$ \\
\midrule
L & $L(x)\,(x-1)=x^{11}-x^{9}-x^{8}+x^{3}+x^{2}-1$ & trace $0$ (commutator); $\Mah=1.17628\in(1,\gold)$; charge group $\bot$ \\
\midrule
M & $\gold\tens\gold$ (on-circle regression) & charpoly $(x{+}1)^2(x^2{-}3x{+}1)$; $\Mah=\gold^{2}$ (tropical), \emph{not} $\gold^{4}$ \\
\bottomrule
\end{tabular}
\end{center}

\section*{References}
\small
\begin{enumerate}[label={[\arabic*]},leftmargin=2.2em,nosep]
  \item L.~Kronecker, \emph{Zwei S\"atze \"uber Gleichungen mit ganzzahligen Coefficienten},
        J.~Reine Angew.~Math.\ \textbf{53} (1857), 173--175.
  \item D.~H.~Lehmer, \emph{Factorization of certain cyclotomic functions},
        Ann.~of Math.\ \textbf{34} (1933), 461--479.
  \item R.~Salem, \emph{Power series with integral coefficients},
        Duke Math.~J.\ \textbf{12} (1945), 153--172.
  \item K.~Mahler, \emph{On some inequalities for polynomials in several variables},
        J.~London Math.~Soc.\ \textbf{37} (1962), 341--344.
  \item C.~J.~Smyth, \emph{On the product of the conjugates outside the unit circle of an
        algebraic integer}, Bull.~London Math.~Soc.\ \textbf{3} (1971), 169--175.
  \item D.~W.~Boyd, \emph{Speculations concerning the range of Mahler's measure},
        Canad.~Math.~Bull.\ \textbf{24} (1981), 453--469.
  \item E.~Dobrowolski, \emph{On a question of Lehmer and the number of irreducible factors of
        a polynomial}, Acta Arith.\ \textbf{34} (1979), 391--401.
  \item K.~Shoda, \emph{Einige S\"atze \"uber Matrizen}, Japan.\ J.\ Math.\ \textbf{13} (1936),
        361--365.
  \item A.~A.~Albert and B.~Muckenhoupt, \emph{On matrices of trace zero},
        Michigan Math.~J.\ \textbf{4} (1957), 1--3.
  \item T.~J.~Laffey and R.~Reams, \emph{Integral similarity and commutators of integral
        matrices}, Linear Algebra Appl.\ \textbf{197/198} (1994), 671--689.
  \item M.~J.~Bertin et al., \emph{Pisot and Salem Numbers}, Birkh\"auser, 1992.
  \item \emph{The $\Zmod5$ Case of the No--Salem Dichotomy}, companion note,
        math-research-pipelines (2026); the source of \S\ref{sec:z5}.
\end{enumerate}
\smallskip
\noindent\textit{Citation caveat:} these references are standard attributions for the results
invoked (Kronecker's theorem, Lehmer's problem, Salem/Pisot numbers, the Mahler measure,
Smyth's non-reciprocal bound, and the Shoda / Albert--Muckenhoupt / Laffey--Reams commutator
results), supplied from memory. Volume, page, and year details have \emph{not} been
machine-verified and should be checked against the primary sources before any external use.

\end{document}
